\section{LectureNotes.cls}\label{sec:ref-lecturenotes}

The lecture class: handouts and their projector twin, from one source.
Based on \file{extarticle} for the large sizes -- 14\,pt handout, 17\,pt
projector. Loads Exercises, ProblemMeta, Workspace, Typefaces (lining),
MathStuff, Titling, Hyperlinks, CourseBoxes and InstructorNotes -- the full
stack. Workflow: section~\ref{sec:lectures}.

\subsection*{Class options}

\begin{center}
\begin{tabularx}{\linewidth}{@{}l Y@{}}
\texttt{projector} & the in-room copy: 17\,pt, tighter top margin, no
  marginpar column, no page number, and the text block sitting left of
  centre (same width, shifted -- see below).
  Normally you never write it -- the build's \env{projector} view injects
  \cs{ProjectorView} and the class switches by itself.\\
\emph{(other)} & forwarded to \file{extarticle}.\\
\end{tabularx}
\end{center}

\subsection*{Commands and hooks}

\begin{center}
\begin{tabularx}{\linewidth}{@{}l Y@{}}
\cs{NotesProblemLabelWord} & the noun imported problems render as
  (default \emph{Example}); redefine per \emph{text}, in
  \file{.textbook/textbook.tex}. Routed into \cs{ProblemLabelWord} so a
  shared text layer cannot disturb Assessment's fixed ``Q''.\\
\cs{lecturetoc} & the lecture's table of contents. Document-called, not
  class-emitted: a short lecture wants none, and only you know which those
  are. Put it straight after \cs{booksection} (lectures have no title
  block). See ``The contents list'' below.\\
\cs{CourseBoxTocEntry} & CourseBoxes' contents hook, rendered here. Renew it
  to change what an entry says.\\
\cs{booksection}, \cs{lecturetitle} & the lecture's title and, for
  \cs{booksection}, the section of the book it covers (with an optional
  \texttt{link=} to that section's page). Supplied by TextRef
  (section~\ref{sec:ref-textref}); the class renders it as the running
  head.\\
\end{tabularx}
\end{center}

\subsection*{Presentation choices}

Imported problems render with a run-in bold heading, unnumbered by default:
\textbf{Example.}\ -- plus the literal \texttt{number=} and parenthetical
\texttt{name=} if the import supplied them (``Example 2.1.8
(differentiating a polynomial).'').\ The problem counter still steps invisibly, so cross-references
stay sound. The \env{problemgroup} stem takes the same unnumbered form.
Points keys are ignored -- notes are ungraded.

Other deliberate class behaviour:

\begin{itemize}
  \item \cs{subsectionbreak} is \cs{clearpage}: every subsection starts a
    fresh page (a lecture subsection is a screenful/boardful).
  \item Only \cs{section} is numbered (\env{secnumdepth} 1); the TOC shows
    three levels.
  \item Paragraph spacing is airier than the parskip default
    (\texttt{skip=1em}) -- room to annotate, legible from the back row.
  \item The running head is the lecture's \cs{booksection} coordinate
    (linked if the book is online), on every page \emph{including} the
    projector copy -- a scrolled
    document is where persistent orientation is worth most. It carries a
    rule, which disappears along with the head if the lecture declares no
    coordinate. Page \emph{numbers} are the opposite case: the handout has a
    centred one, the projector copy none (it is never collated, so a number
    is noise); the projector foot is simply empty.
  \item The projector text block is shifted left of centre
    (\texttt{hmarginratio=\{3:7\}}, so the left margin is 60\% of the centred
    one and the rest falls to the right). Its \emph{width} is the same measure
    as everywhere else -- this moves the block, it does not reflow it. The
    reason is the room: presenting from a tablet puts the app's tool bar down
    the right edge of the screen, over the projected text.
  \item Headings and links stay black: lectures sit with assessments on
    the ``clear, not fancy'' side; colour belongs to CourseDocument.
\end{itemize}

\subsection*{The contents list}

A lecture's contents list is \emph{not} mostly sections -- most lectures
have none. It is mostly the things that actually structure the hour, and it
is assembled from three sources:

\begin{itemize}
  \item every CourseBoxes \textbf{text result} -- definitions, theorems,
    facts, examples, and whatever else the text layer declares -- via
    \cs{CourseBoxTocEntry}. \textbf{Asides never appear}, and not by a list
    of exceptions: they are a different factory, so the exclusion cannot
    drift when a text layer adds a name.
  \item every \textbf{top-level imported problem}, through the same
    renderer, so ``Example 2.1.8'' and ``Definition 2.1.1'' cannot end up
    formatted differently. Parts inside a \env{problemgroup} and check-ins
    stay out for free -- neither reaches the heading hook.
  \item a \env{problemgroup}, \emph{if} it was given a \texttt{name=}. The
    group's stem is a sentence and reads badly in a list, so the name is what
    gets listed; without one the group stays out, since a bare ``Example''
    would be indistinguishable from any other.
  \item \cs{section}/\cs{subsection}, when a lecture uses them.
\end{itemize}

Entries read \emph{Kind number (name)}, dropping whatever is absent:
``Definition 2.1.1 (Derivative)'', ``Example 2.1.5'', ``Theorem (Product,
Quotient, and Chain Rules)''.

Two deliberate presentation choices:

\begin{itemize}
  \item \textbf{Flat, with the typeface carrying the hierarchy.}
    Indentation implies structure; in a sectionless lecture there is none to
    imply, and indenting every result would push the list right for nothing.
    So all levels sit at one margin and sans-bold (headings, echoing the
    page) separates from serif roman (results).
  \item \textbf{Page numbers on the handout only.} The projector copy prints
    no page numbers anywhere, so a page column there would point at numbers
    the reader cannot see. Its entries are hyperlinks instead -- black, so
    they read as plain text and behave as taps.
\end{itemize}
